\subsection{外部化指数、失真阈值与模性拓扑的统一命题链}\label{subsec:conclusion-externalization-distortion-modularity-topology}
本小节把外部化指数、率失真硬下界、\(\varepsilon\)-sound 稳定性、异常放大动力学以及 Lee--Yang 的 \(S_4\) 模性接口压缩为一组可独立调用的命题，并与前文已证结构保持同一依赖链。

\begin{theorem}[外部化指数控制的最坏读出时间下界]\label{thm:conclusion-externalization-index-readout-time-lower-bound}
固定分辨率 \(m\)，令
\[
\Fold_m:\Omega_m\to X_m,\qquad
d_m(x):=\card{\Fold_m^{-1}(x)},\qquad
D_m:=\max_{x\in X_m}d_m(x).
\]
考虑任意顺序读出--解码协议：在给定 \(x\in X_m\) 条件下，协议逐步读取字母表 \(\mathcal A\)（\(|\mathcal A|\ge 2\)）并在停止时输出 \(\widehat\omega\in \Fold_m^{-1}(x)\)，且要求在每条纤维上实现零误差单值复原。设 \(\mathbb E[T(x)]\) 为纤维 \(\Fold_m^{-1}(x)\) 上均匀先验下的停止时间期望，则
\[
\sup_{x\in X_m}\mathbb E[T(x)]
\ \ge\
\frac{\log D_m}{\log|\mathcal A|}.
\]
此外，对折叠诱导的条件期望 \(E_m\)（命题 \ref{prop:fold-condexp-index-maxfiber}），有
\[
\Ind(E_m)=D_m,
\]
从而同一界可写为
\[
\sup_{x\in X_m}\mathbb E[T(x)]
\ \ge\
\frac{\log \Ind(E_m)}{\log|\mathcal A|}.
\]
\end{theorem}
\begin{proof}
取 \(x_\star\in X_m\) 使 \(d_m(x_\star)=D_m\)。在纤维 \(\Fold_m^{-1}(x_\star)\) 上，零误差顺序读出对应一棵 \(|\mathcal A|\)-叉前缀树，其叶长记为 \(\ell_1,\dots,\ell_{D_m}\)。由 Kraft 不等式，
\[
\sum_{i=1}^{D_m}|\mathcal A|^{-\ell_i}\le 1.
\]
在均匀先验下，
\[
\mathbb E[T(x_\star)]=\frac1{D_m}\sum_{i=1}^{D_m}\ell_i.
\]
由上式与对数凸性（或等价的前缀码平均长度下界）得
\[
\frac1{D_m}\sum_{i=1}^{D_m}\ell_i\ge \log_{|\mathcal A|}D_m
=\frac{\log D_m}{\log|\mathcal A|}.
\]
故 \(\sup_x\mathbb E[T(x)]\ge \mathbb E[T(x_\star)]\ge \log D_m/\log|\mathcal A|\)。
再由命题 \ref{prop:fold-condexp-index-maxfiber} 的 \(\Ind(E_m)=D_m\) 即得指数表述。证毕。
\end{proof}

\begin{corollary}[面积律饱和下的边界读出时间下界]\label{cor:conclusion-area-law-page-time-lower-bound}
若外部接口满足面积律饱和
\[
S_{\mathrm{BH}}=\frac{A}{4},
\]
并按本稿容量口径满足
\[
S_{\mathrm{BH}}=\log\Ind(E_m),
\]
则任意 \(|\mathcal A|\)-元顺序边界读出的最坏期望时间满足
\[
\sup_{x}\mathbb E[T(x)]
\ \ge\
\frac{A}{4\log|\mathcal A|}.
\]
特别地，当 \(|\mathcal A|=2\) 时，
\[
\sup_x\mathbb E[T(x)]\ge \frac{A}{4}.
\]
\end{corollary}
\begin{proof}
由命题 \ref{prop:dephys-horizon-fiber-capacity-area-law} 与定理 \ref{thm:conclusion-externalization-index-readout-time-lower-bound} 直接合并。证毕。
\end{proof}

\begin{theorem}[Fibonacci 稳定型输出的不可消除失真阈值]\label{thm:conclusion-fibonacci-hard-distortion-threshold}
令 \(\omega\sim\mathrm{Unif}(\{0,1\}^m)\)，\(X=\Fold_m(\omega)\)。
对任意（可随机化）解码器 \(\widehat\omega=\mathsf{Dec}(X)\)，定义归一化期望汉明失真
\[
\delta:=\frac1m\,\mathbb E[d_H(\omega,\widehat\omega)]\in[0,1/2].
\]
记二元熵函数
\[
h_2(t):=-t\log_2 t-(1-t)\log_2(1-t).
\]
则必有
\[
h_2(\delta)\ \ge\ 1-\frac1m\log_2|X_m|.
\]
若进一步 \(|X_m|=F_{m+2}\)，并以 \(\delta_m^\star\in(0,1/2)\) 定义
\[
h_2(\delta_m^\star)=1-\frac1m\log_2F_{m+2},
\]
则对任意解码器有 \(\delta\ge\delta_m^\star\)。并且当 \(m\to\infty\) 时，
\[
\delta_m^\star\to\delta_\infty^\star,\qquad
h_2(\delta_\infty^\star)=1-\log_2\varphi,
\]
数值上 \(\delta_\infty^\star\approx 0.054632351165136\)。
\end{theorem}
\begin{proof}
由数据处理不等式，
\[
I(\omega;\widehat\omega)\le I(\omega;X)\le H(X)\le \log_2|X_m|.
\]
另一方面，二元均匀源在平均汉明失真 \(\delta\) 下的率失真函数为 \(R(\delta)=1-h_2(\delta)\)（\(\delta\in[0,1/2]\)），故
\[
I(\omega;\widehat\omega)\ge m(1-h_2(\delta)).
\]
合并即得
\[
m(1-h_2(\delta))\le \log_2|X_m|,
\]
等价于所述不等式。
当 \(|X_m|=F_{m+2}\) 时，\(h_2\) 在 \([0,1/2]\) 上严格单调，故 \(\delta_m^\star\) 唯一并给出 \(\delta\ge\delta_m^\star\)。
再由 \(F_{m+2}\sim \varphi^m/\sqrt5\) 得
\(
\frac1m\log_2F_{m+2}\to\log_2\varphi
\)，从而得到极限方程与数值。证毕。
\end{proof}

\begin{corollary}[侧信息预算与圆维下界]\label{cor:conclusion-sideinfo-distortion-cdim-lower-bound}
在定理 \ref{thm:conclusion-fibonacci-hard-distortion-threshold} 的设定下，若引入侧信息 \(S\)（比特预算 \(L:=H(S)\)）并用 \(\widehat\omega=\mathsf{Dec}(X,S)\) 达到失真 \(\delta\)，则
\[
L\ \ge\ m(1-h_2(\delta))-\log_2|X_m|.
\]
若侧信息由相位圆离散化实现：设可见圆轴数为 \(\cdim(G)\)，每轴分辨率为 \(Q\)（每轴至多 \(\log_2 Q\) 比特），则必有
\[
\cdim(G)\ \ge\
\frac{m(1-h_2(\delta))-\log_2|X_m|}{\log_2 Q}.
\]
在 \(|X_m|=F_{m+2}\) 且 \(\delta<\delta_m^\star\) 时，右端为正；特别地，若 \(Q\) 为常数则 \(\cdim(G)=\Omega(m)\)，若 \(Q=\mathrm{poly}(m)\) 则 \(\cdim(G)=\Omega(m/\log m)\)。
\end{corollary}
\begin{proof}
与定理 \ref{thm:conclusion-fibonacci-hard-distortion-threshold} 同理，
\[
I(\omega;\widehat\omega)\le I(\omega;X,S)\le H(X)+H(S)\le \log_2|X_m|+L.
\]
而率失真下界仍给出
\(
I(\omega;\widehat\omega)\ge m(1-h_2(\delta))
\)。
故
\[
L\ge m(1-h_2(\delta))-\log_2|X_m|.
\]
若每个圆轴最多承载 \(\log_2Q\) 比特，则 \(H(S)\le \cdim(G)\log_2Q\)，代回即得圆维下界。最后两条增长律由 \(|X_m|=F_{m+2}\sim\varphi^m\) 的线性主阶直接给出。证毕。
\end{proof}

\begin{corollary}[Gödel 整数封装的超线性不可压缩性]\label{cor:conclusion-godel-sideinfo-superlinear-overhead}
设侧信息最终以 prime-register 向量 \(v\) 封装为单个 Gödel 整数 \(N(v)\)，其向量长度记为 \(L(v)\)。
则由定理 \ref{thm:xi-godel-integerization-log-overhead}，存在绝对常数 \(c_0,c_1>0\) 使
\[
\log_2N(v)\ \ge\ c_0\,L(v)\log_2L(v)-c_1L(v).
\]
特别地，在 \(|X_m|=F_{m+2}\) 且 \(\delta<\delta_\infty^\star\) 时，由推论 \ref{cor:conclusion-sideinfo-distortion-cdim-lower-bound} 得 \(L(v)=\Omega(m)\)，从而
\[
\log_2N(v)=\Omega(m\log m).
\]
\end{corollary}
\begin{proof}
第一式即定理 \ref{thm:xi-godel-integerization-log-overhead} 的量纲表述。
第二式由 \(L(v)\ge L=\Omega(m)\) 代入第一式即得。证毕。
\end{proof}

\begin{theorem}[面积律预算下标量 Gödel 外置的结构不相容性]\label{thm:conclusion-area-law-scalar-godel-incompatibility}
设边界可见自由度规模为 \(b\)，并假设存在与 \(b\) 无关的常数 \(C>0\) 使得边界实现层满足严格面积律预算
\[
\log_2\dim(\mathcal H_{\partial})\ \le\ C b.
\]
若试图以互异素数 \(q_1,\dots,q_b\) 将 \(\{0,1\}^b\) 标量 Gödel 化为
\(
g(\varepsilon):=\prod_{i=1}^b q_i^{\varepsilon_i}
\)
并以最坏情形 \(L_{\max}:=\max_{\varepsilon}g(\varepsilon)=\prod_{i=1}^b q_i\) 的位长 \(\log_2 L_{\max}\) 作为实现层资源度量，则必有
\[
\log_2 L_{\max}\ \ge\ \log_2(p_b\#)\ =\ \Omega(b\log b),
\]
其中 \(p_b\#:=\prod_{i=1}^b p_i\) 为 primorial（\(p_i\) 为第 \(i\) 个素数）。
因此当 \(b\to\infty\) 时，该标量外置的最坏资源增长严格超出任何线性面积律预算。
\end{theorem}
\begin{proof}
由素数序性质 \(q_i\ge p_i\) 得
\(
L_{\max}=\prod_{i=1}^b q_i\ge \prod_{i=1}^b p_i=p_b\#
\)，从而 \(\log_2 L_{\max}\ge \log_2(p_b\#)\)。
又由 \(p_i\ge i+1\) 得 \(p_b\#\ge (b+1)!\)，故 \(\log_2(p_b\#)\ge \log_2((b+1)!)=\Omega(b\log b)\)。
结合线性面积律预算 \(\log_2\dim(\mathcal H_{\partial})\le Cb\) 即得结论。证毕。
\end{proof}

\begin{theorem}[\texorpdfstring{\(\varepsilon\)}{eps}-sound 预算的 Wasserstein--Stieltjes 解析稳定性]\label{thm:conclusion-epssound-w1-stieltjes-stability}
固定接口对象 \(\mathcal I\) 并沿定义 \ref{def:pom-eps-sound-rewrite} 记
\[
w\Longrightarrow_\varepsilon w'.
\]
设同一读出通道给出的随机变量 \(A_w,A_{w'}\) 在同一概率空间上取值于 \([1,B]\)，且在好事件上满足 \(A_w=A_{w'}\)。
记其分布为 \(\mu_w,\mu_{w'}\)。
则
\[
W_1(\mu_w,\mu_{w'})\le 2B\varepsilon.
\]
进一步，对任意 \(t>0\) 的 Stieltjes 读出
\[
S_{\mu}(t):=\int\frac{1}{t+a}\,\dd\mu(a),
\]
有显式稳定界
\[
\bigl|S_{\mu_w}(t)-S_{\mu_{w'}}(t)\bigr|
\le
\frac{2B\varepsilon}{(t+1)^2}.
\]
\end{theorem}
\begin{proof}
由定义 \ref{def:pom-eps-sound-rewrite}，存在坏事件 \(E\) 满足 \(\mathbb P(E)\le\varepsilon\)，且在 \(E^c\) 上 \(A_w=A_{w'}\)。
故
\[
\mathbb E|A_w-A_{w'}|
=\mathbb E\!\bigl[|A_w-A_{w'}|\mathbf 1_E\bigr]
\le 2B\,\mathbb P(E)\le 2B\varepsilon.
\]
由 Wasserstein 距离的耦合定义，\(W_1(\mu_w,\mu_{w'})\le \mathbb E|A_w-A_{w'}|\)，从而得第一式。
再由命题 \ref{prop:xi-stieltjes-w1-lipschitz}，
\[
\bigl|S_{\mu_w}(t)-S_{\mu_{w'}}(t)\bigr|
\le \frac{1}{(t+1)^2}W_1(\mu_w,\mu_{w'}),
\]
代入即可。证毕。
\end{proof}

\begin{theorem}[异常放大在环面上的扩张动力学与熵律]\label{thm:conclusion-anomaly-torus-expanding-entropy}
设 lift 角色空间给出实向量空间
\[
A_G:=\RR^{\widehat G\setminus\{\chi_0\}},
\]
取满秩格 \(\Lambda\subset A_G\)，并记环面
\[
\mathbb T_G:=A_G/\Lambda,\qquad d:=\dim\mathbb T_G.
\]
对任意整数 \(q\ge 2\)，定义
\[
F_q:\mathbb T_G\to\mathbb T_G,\qquad F_q([a])=[qa].
\]
则：
\begin{enumerate}
  \item \(F_q\) 为 \(C^\infty\) 扩张覆叠映射，度数为 \(q^d\)；
  \item 非游荡集等于 \(\mathbb T_G\) 本身；
  \item 拓扑熵满足
  \[
  h_{\mathrm{top}}(F_q)=d\log q.
  \]
\end{enumerate}
此外，若异常类满足矩放大律
\[
[\mathrm{Anom}(\mathrm{MOM}_q\circ w)]=q\,[\mathrm{Anom}(w)]
\]
（命题 \ref{prop:pom-anom-moment-amplification}），则其在环面坐标上的动力学由 \(F_q\) 给出。
\end{theorem}
\begin{proof}
线性映射 \(a\mapsto qa\) 在商空间 \(\mathbb T_G\) 上诱导 \(F_q\)，其雅可比行为为常数 \(qI_d\)，故处处扩张且覆盖度为 \(|\det(qI_d)|=q^d\)。
扩张环面端同态的非游荡集为全空间，且拓扑熵等于对数体积增长率 \(d\log q\)。
最后一条由命题 \ref{prop:pom-anom-moment-amplification} 的线性放大公式与环面商映射兼容性直接得到。证毕。
\end{proof}

\begin{theorem}[Lee--Yang 的 \texorpdfstring{$S_4$}{S4} 专化、奇二维 Artin 表示与权一模性]\label{thm:conclusion-leyang-s4-odd-artin-weight-one}
令
\[
P_{\mathrm{LY}}(y)=y^4-3y^3+4y^2-2y+1.
\]
取 \(y_0\in\QQ\setminus\{0,1\}\) 且 \(P_{\mathrm{LY}}(y_0)\neq 0\)，并设对应四次特化的分裂域满足
\[
\Gal(L_{y_0}/\QQ)\cong S_4.
\]
记
\[
K_{y_0}:=\QQ\!\left(\sqrt{-y_0(y_0-1)P_{\mathrm{LY}}(y_0)}\right).
\]
若 \(K_{y_0}\) 为虚二次域，则由 \(S_4\) 的二维不可约表示诱导的 Artin 表示 \(\rho_{y_0}\) 为奇表示，并存在权 \(1\) 本原尖点新形式 \(f_{y_0}\) 使
\[
L(s,f_{y_0})=L(s,\rho_{y_0}).
\]
并且可取无穷多 \(y_0>1\) 使上述结论同时成立。
\end{theorem}
\begin{proof}
由定理 \ref{thm:xi-time-part9c-leyang-cover-s4-specialization} 与定理 \ref{thm:xi-time-part9c-leyang-s4-specialization-artin-weight1}，\(S_4\) 正则专化可诱导二维 Artin 表示，并在奇性条件下对应权一模形式。
当 \(-y_0(y_0-1)P_{\mathrm{LY}}(y_0)<0\) 时，\(K_{y_0}\) 为虚二次域，从而满足奇性判别。
另一方面，\(P_{\mathrm{LY}}(y)>0\) 在充分大 \(y>1\) 上成立，故上述负性条件在无穷多有理点成立；再由 Hilbert 不可约性可取无穷多 \(y_0\) 保持 \(S_4\) 专化。证毕。
\end{proof}

\begin{corollary}[折叠 \texorpdfstring{\(\chi\)}{chi}-twisted \texorpdfstring{\(\zeta\)}{zeta} 的条件模性闭包]\label{cor:conclusion-chi-twisted-zeta-conditional-modularity}
设存在二维 Artin 表示 \(\rho\) 满足：除有限素数集外，折叠 \(\chi\)-twisted \(\zeta\) 的局部 Euler 因子与 \(L(s,\rho)\) 局部因子一致。
若 \(\rho\) 满足定理 \ref{thm:conclusion-leyang-s4-odd-artin-weight-one} 的奇性与可解像条件，则存在权一新形式 \(f\) 使
\[
L(s,f)=L(s,\rho),
\]
并且该 \(\chi\)-twisted \(\zeta\) 的原始因子在除有限素数外与 \(L(s,f)\) 一致，故继承相应的解析延拓与函数方程结构。
\end{corollary}
\begin{proof}
由定理 \ref{thm:conclusion-leyang-s4-odd-artin-weight-one}，\(L(s,\rho)\) 来自权一模形式。
局部因子在有限外一致即给出 Euler 乘积的有限处同调修正，故原始部分同一并继承权一模形式的解析性质。证毕。
\end{proof}

\begin{conjecture}[Fold--Jones 谱拓扑同构]\label{conj:conclusion-fold-jones-spectral-topological-equivalence}
存在自然 \(2\)-函子
\[
\mathcal J:\mathbf{PW}^{(\delta)}\longrightarrow \mathrm{TL}_{\delta},
\qquad
\delta=\sqrt{\Ind(E_m)},
\]
其中 \(\mathbf{PW}^{(\delta)}\) 为折叠 witness-block 可计算切片，\(\mathrm{TL}_{\delta}\) 为 loop 参数 \(\delta\) 的 Temperley--Lieb 平面代数。
对任意由投影词编码的闭辫 \(L\)，其 Jones 评估与原始 \(\zeta\)-行列式读出满足同值对应；并且 Lee--Yang 参数 \(y\) 的解析延拓在 \((q,y)\)-平面上诱导受控分歧集合，与局部 Artin 因子匹配保持同构。
该猜想与猜想 \ref{conj:pom-fold-jones-tower-symtft-doubled-fib}、猜想 \ref{conj:pom-fold-leyang-jones-boundary-value} 在接口层一致。
最小可行验证路径为：在固定 \(m\) 的 witness-block 上显式构造 Jones 投影，核验 Temperley--Lieb 关系，并比较 Markov 迹与 \(\zeta\)-读出分母多项式。
\end{conjecture}

